%! Author = lili
%! Date = 6/9/26

%V Seiten in Genetik von Knust / Janning:
% 68-75: Analyse von Erbgängen -  Multiple Allelie
%% https://link.springer.com/chapter/10.1007/978-3-662-70442-4_7#Sec12
% 35-41: Meiose - Die Meiose genetisch gesehen ( Unterschiede in der zytologischen und genetischen Betrachtung
%der Meiose + Wann werden die Zellen während der Meiose haploid? )
% 84-104: Genkartierung - Wie kann man genetische Kopplung erkennen?, Testkreuzung zur Interpretation der
% Kopplungsverhältnisse, Statistik: Stimmen Hypothese und Experiment überein?, Dreifaktorenkreuzungen, Tetradenanalyse
% 148-151: Struktur und Funktion der DNA - Genkonversion

% Bis F 9 Wdh.
\section{Sequenzbasierte Molekulare Marker}
\defin{marker}
\begin{fr}
    \textbf{Wie kann man DNA-Varianten zwischen Familienmitgliedern oder möglichen Kandidaten nachweisen?}
    \begin{enumerate}
        \item Man nimmt DNA-Proben von Personen, z. B. Eltern, Kindern oder möglichen Verwandten/Kandidaten.
        \item Mit der PCR (Polymerase-Kettenreaktion) wird ein bestimmter Abschnitt der DNA vervielfältigt.
        \item Die PCR-Produkte (Amplicons) werden in einem Agarosegel nach ihrer Größe getrennt.
        \item Man schaut, welche PCR-Produkte entstanden sind.
    \end{enumerate}
    Die Unterschiede entstehen durch Veränderungen in der DNA-Sequenz.
\end{fr}
\zitat{Bei molekularen Markern kann es sich um einfache, hintereinanderliegende, also
\textbf{in Tandem angeordnete Di-, Tri-, Tetra‐Penta- oder Hexanukleotide} handeln, die verstreut im Genom lokalisiert sind und die Gruppe der \glspl{ssr} bilden.}{\cite{Knust2025_Analyse_von_Erbgängen}}
\paragraph{Marker können sein:}
\begin{itemize}
    \item Längenvariationen
    \item Sequenzvariationen
\end{itemize}

\begin{itemize}
    \item \gls{str} bzw. kurze DNA-Wiederholungen (\gls{mikrosatellit} / \gls{ssr}): Die Sequenz selbst ist gleich,
    aber die Anzahl der Wiederholungen ist unterschiedlich.
    \begin{itemize}
        \item Jeder Mensch besitzt eine bestimmte Kombination solcher Wiederholungszahlen.
        \item Diese Varianten können von Eltern an Kinder weitergegeben werden.
    \end{itemize}
    \item \gls{rflp}
    \item \gls{snp}
    \item \gls{vntr}
    \begin{itemize}
        \item Im Genom gibt es verschiedene Stellen, an denen man VNTRs findet. Diese Loci werden nach Genen benannt, neben denen sie auftreten. Im Gegensatz zu Mikrosatelliten (STR) sind sie längere repetitive Abschnitte.
    \end{itemize}
    \item \gls{rapd}
\end{itemize}


\begin{bsp}
    \textbf{Vaterschaftstest mit molekularen Markern}
\end{bsp}

\section{Chromatin - DNA Verpackung im Zellkern}

\section{Allelfrequenzen in Populationen}

\section{Genkonversion}